\documentclass[10pt]{article}
\usepackage[paperwidth=42cm, paperheight=29.7cm, margin=0pt]{geometry}
\usepackage{fontspec}
\usepackage{microtype}
\usepackage[brazil]{babel}
\usepackage{amsmath}
\usepackage{amssymb}
\usepackage{tikz}
\usetikzlibrary{calc, positioning, arrows.meta, shapes.geometric, decorations.pathreplacing, fit, backgrounds}

% === TIPOGRAFIA (Dark Academic) ===
\setmainfont{Athelas}[Ligatures=TeX]
\setsansfont{Avenir Next}
\newfontfamily\headerfont{Avenir Next}[LetterSpace=6]
\newfontfamily\titlefont{Avenir Next Condensed}[LetterSpace=10]

% === PALETA DARK ===
\definecolor{darkbg}{RGB}{15, 23, 42}       % Slate 900
\definecolor{cardbg}{RGB}{30, 41, 59}       % Slate 800
\definecolor{cardborder}{RGB}{51, 65, 85}   % Slate 700
\definecolor{lighttext}{RGB}{226, 232, 240} % Slate 200
\definecolor{dimtext}{RGB}{148, 163, 184}   % Slate 400
\definecolor{cyan}{RGB}{6, 182, 212}        % Ciano acento
\definecolor{rose}{RGB}{244, 63, 94}        % Rosa acento
\definecolor{amber}{RGB}{245, 158, 11}      % Amber acento
\definecolor{emerald}{RGB}{16, 185, 129}    % Emerald acento

\usepackage{xcolor}
\pagecolor{darkbg}
\color{lighttext}
\pagestyle{empty}

% === MACROS DE CAIXA DARK ===
\newcommand{\darkbox}[3]{% {acento}{título}{conteúdo}
    \begin{tikzpicture}
        \node[rectangle, rounded corners=6pt, fill=cardbg, draw=cardborder,
              line width=0.6pt, inner sep=14pt,
              text width=\linewidth-28pt, align=left] (box) {%
            {\headerfont\normalsize\color{#1}\MakeUppercase{#2}}\\[6pt]
            {\small\color{lighttext}#3}
        };
        % Acento lateral
        \fill[#1, rounded corners=2pt] ([xshift=-0.5pt]box.north west) rectangle ([xshift=3pt, yshift=0pt]box.south west);
    \end{tikzpicture}\par\vspace{0.4cm}
}

\newcommand{\titem}[1]{%
    \noindent\hangindent=1.5em\hangafter=1%
    {\color{cyan}\footnotesize$\blacktriangleright$}\hspace{0.5em}#1\par\vspace{3pt}%
}

\begin{document}

% === FUNDO COM GRID (overlay) ===
\begin{tikzpicture}[remember picture, overlay]
    % Grid sutil
    \foreach \x in {0,1,...,42} {
        \draw[cardborder!15] (\x, 0) -- (\x, -29.7);
    }
    \foreach \y in {0,-1,...,-29.7} {
        \draw[cardborder!15] (0, \y) -- (42, \y);
    }
\end{tikzpicture}

% ==================================================================
% HEADER
% ==================================================================
\noindent
\begin{tikzpicture}[remember picture, overlay]
    \fill[cardbg] (current page.north west) rectangle ([yshift=-4cm]current page.north east);
    \fill[cyan] ([yshift=-4cm]current page.north west) rectangle ([yshift=-4.15cm]current page.north east);
\end{tikzpicture}

\noindent
\begin{minipage}[c][\dimexpr4cm\relax][c]{\textwidth}
    \hspace{2.5cm}%
    \begin{minipage}[c]{0.65\textwidth}
        {\titlefont\fontsize{40}{44}\selectfont\color{white}\MakeUppercase{Oximetria de Pulso}}\\[4pt]
        {\headerfont\color{dimtext}Monitorização Contínua e Não-Invasiva da Saturação Arterial}
    \end{minipage}\hfill
    \begin{minipage}[c]{0.25\textwidth}
        \begin{flushright}
            {\headerfont\color{cyan}\Large ONE PAGER ICU}\\[3pt]
            {\small\color{dimtext}\textit{Nick Mark, MD}}\\[1pt]
            {\footnotesize\color{cardborder}onepagericu.com}
        \end{flushright}
    \end{minipage}%
    \hspace{2.5cm}
\end{minipage}

\vspace{1cm}

% ==================================================================
% CORPO PRINCIPAL
% ==================================================================
\noindent\hspace{2.5cm}%
\begin{minipage}[t]{\dimexpr\textwidth-5cm\relax}

% --- PRINCÍPIO ---
\darkbox{cyan}{Princípio Fisiológico}{%
    A oximetria de pulso constitui a medição contínua da saturação de O\textsubscript{2} da hemoglobina por espectrofotometria. O sensor emite luz \textbf{\color{rose}vermelha} (660\,nm) e \textbf{\color{amber}infravermelha} (940\,nm) através do leito capilar. A absorção diferencial entre hemoglobina \textbf{oxigenada} (OxyHb) e \textbf{desoxigenada} (DeoxyHb) é analisada usando o componente pulsátil (AC) do sinal.
}

\vspace{0.2cm}

% --- 2 COLUNAS: SENSOR + GRÁFICO ---
\noindent
\begin{minipage}[t]{0.48\textwidth}
    \darkbox{amber}{Mecânica do Sensor}{%
        \begin{center}
        \begin{tikzpicture}[scale=1.0, font=\sffamily\small]
            \draw[thick, dimtext, fill=cardbg!50] (-2.2, -0.7) .. controls (1, -0.7) and (1.5, -0.5) .. (1.8, 0) .. controls (1.5, 0.5) and (1, 0.7) .. (-2.2, 0.7);
            \draw[cardborder] (0.8, 0.5) .. controls (1.5, 0.5) and (1.8, 0.3) .. (1.8, 0) .. controls (1.8, -0.3) and (1.5, -0.5) .. (0.8, -0.5);
            \fill[cardborder] (-0.8, 1.0) rectangle (1.0, 1.35);
            \fill[cardborder] (-0.8, -1.0) rectangle (1.0, -1.35);
            \node[lighttext, font=\tiny\headerfont] at (0.1, 1.17) {LED};
            \node[lighttext, font=\tiny\headerfont] at (0.1, -1.17) {SENSOR};
            \draw[dashed, rose, -Latex, thick] (0.0, 0.95) -- (0.0, -0.95);
            \draw[dashed, amber, -Latex, thick] (0.4, 0.95) -- (0.4, -0.95);
            \node[right, rose, font=\tiny] at (0.0, 0.5) {660nm};
            \node[right, amber, font=\tiny] at (0.4, -0.5) {940nm};
        \end{tikzpicture}
        \end{center}
        O componente \textbf{AC} (pulsátil) isola sangue arterial.\\
        O componente \textbf{DC} (constante) filtra tecidos e sangue venoso.
    }
\end{minipage}\hfill
\begin{minipage}[t]{0.48\textwidth}
    \darkbox{emerald}{Espectro de Absorção}{%
        \begin{center}
        \begin{tikzpicture}[scale=0.80, font=\sffamily\small]
            \draw[->, dimtext] (0,0) -- (6.2,0) node[right, scale=0.7, lighttext] {nm};
            \draw[->, dimtext] (0,0) -- (0,4.2) node[above, scale=0.7, lighttext] {Abs};
            \draw[ultra thick, cyan] (0, 3.5) .. controls (2, 2.8) and (3, 2) .. (5.5, 1) node[right, scale=0.7] {\textbf{OxyHb}};
            \draw[ultra thick, rose] (0, 1) .. controls (1, 2) and (2, 3) .. (3, 3) .. controls (4, 3) and (5, 2.2) .. (5.5, 2.2) node[right, scale=0.7] {\textbf{DeoxyHb}};
            \draw[dotted, rose!50, thick] (0.8, 0) -- (0.8, 4) node[above, scale=0.6, rose] {\textbf{RED}};
            \draw[dotted, amber!50, thick] (4.8, 0) -- (4.8, 4) node[above, scale=0.6, amber] {\textbf{IR}};
        \end{tikzpicture}
        \end{center}
    }
\end{minipage}

\vspace{0.1cm}

% --- 4 COLUNAS: FATORES ---
\noindent
\begin{minipage}[t]{0.24\textwidth}
    \darkbox{rose}{Pigmentação e Esmalte}{%
        \titem{SpO\textsubscript{2} pode \textbf{superestimar} em pele negra --- hipoxemia oculta.}
        \titem{Pacientes negros: \textbf{3$\times$} mais risco.}
        \titem{Esmaltes escuros (azul, preto) interferem.}
    }
\end{minipage}\hfill
\begin{minipage}[t]{0.24\textwidth}
    \darkbox{cyan}{Hemoglobinas Anormais}{%
        \titem{\textbf{CoHb:} Falsamente normal.}
        \titem{\textbf{MetHb:} Trava em $\sim$85\%.}
        \titem{\textbf{SulfaHb:} Falsa baixa.}
        \titem{Co-oximetria (4+ $\lambda$) resolve.}
    }
\end{minipage}\hfill
\begin{minipage}[t]{0.24\textwidth}
    \darkbox{amber}{Baixo Fluxo e Lag}{%
        \titem{Choque reduz sinal pulsátil.}
        \titem{ECMO/LVAD: inútil.}
        \titem{\textbf{Lag:} 5--15s periférico.}
        \titem{Precisão cai $<$75\%.}
    }
\end{minipage}\hfill
\begin{minipage}[t]{0.24\textwidth}
    \darkbox{emerald}{Corantes e Hiperóxia}{%
        \titem{Azul de metileno: falsa queda.}
        \titem{Hiperóxia não detectada quando SpO\textsubscript{2} = 100\%.}
        \titem{Alvo: SpO\textsubscript{2} $\geq$ 94\%.}
    }
\end{minipage}

\vspace{0.1cm}

% --- ONDAS ---
\noindent
\begin{minipage}[t]{0.30\textwidth}
    \darkbox{cyan}{Onda Pletismográfica}{%
        Componentes \textit{sistólico} e \textit{diastólico} refletem tônus vascular e compliance arterial.\\[4pt]
        \begin{center}
        \begin{tikzpicture}[scale=0.75]
            \fill[cyan!10] (0,0) .. controls (0.2,0) and (0.6,2.2) .. (0.8,2.2) .. controls (1.0,2.2) and (1.2,1) .. (1.5,1) .. controls (1.8,1.4) and (2.5,0) .. (3,0) -- cycle;
            \draw[thick, cyan] (0,0) .. controls (0.2,0) and (0.6,2.2) .. (0.8,2.2) .. controls (1.0,2.2) and (1.2,1) .. (1.5,1) .. controls (1.8,1.4) and (2.5,0) .. (3,0);
            \node[above, cyan, font=\headerfont\tiny] at (0.8, 2.2) {SÍSTOLE};
            \node[right, dimtext, font=\tiny] at (1.5, 1) {Dicrótico};
            \draw[dashed, dimtext!40] (0.8,0) -- (0.8,2.2);
            \draw[dashed, dimtext!40] (1.5,0) -- (1.5,1);
        \end{tikzpicture}
        \end{center}
    }
\end{minipage}\hfill
\begin{minipage}[t]{0.66\textwidth}
    \begin{tikzpicture}
        \node[rectangle, rounded corners=6pt, fill=cardbg, draw=cardborder, line width=0.6pt,
              inner sep=14pt, text width=\linewidth-28pt, align=left] (box) {%
            {\headerfont\normalsize\color{cyan}\MakeUppercase{Padrões de Onda}}\\[10pt]
            \begin{tikzpicture}[scale=0.55, font=\headerfont\scriptsize]
                \node[left, cyan, font=\headerfont\footnotesize\bfseries] at (-0.5, 0.5) {NORMAL};
                \draw[thick, cyan] (0,0) .. controls (0.2,0) and (0.4,1) .. (0.5,1) .. controls (0.6,1) and (0.7,0.4) .. (0.8,0.4) .. controls (0.9,0.8) and (1.5,0) .. (2,0)
                                    (2,0) .. controls (2.2,0) and (2.4,1) .. (2.5,1) .. controls (2.6,1) and (2.7,0.4) .. (2.8,0.4) .. controls (2.9,0.8) and (3.5,0) .. (4,0);
                \begin{scope}[shift={(6,0)}]
                    \node[left, dimtext, font=\headerfont\footnotesize\bfseries] at (-0.5, 0.15) {FRACO};
                    \draw[thick, dimtext] (0,0) .. controls (0.2,0) and (0.4,0.3) .. (0.5,0.3) .. controls (0.6,0.3) and (0.7,0.1) .. (0.8,0.1) .. controls (0.9,0.2) and (1.5,0) .. (2,0)
                                           (2,0) .. controls (2.2,0) and (2.4,0.3) .. (2.5,0.3) .. controls (2.6,0.3) and (2.7,0.1) .. (2.8,0.1) .. controls (2.9,0.2) and (3.5,0) .. (4,0);
                \end{scope}
                \begin{scope}[shift={(12,0)}]
                    \node[left, rose, font=\headerfont\footnotesize\bfseries] at (-0.5, 0.5) {VASOCONSTR.};
                    \draw[thick, rose] (0,0) .. controls (0.2,0) and (0.4,1) .. (0.5,1) .. controls (0.6,1) and (0.7,0.7) .. (0.8,0.7) .. controls (0.9,0.9) and (1.5,0) .. (2,0)
                                        (2,0) .. controls (2.2,0) and (2.4,1) .. (2.5,1) .. controls (2.6,1) and (2.7,0.7) .. (2.8,0.7) .. controls (2.9,0.9) and (3.5,0) .. (4,0);
                \end{scope}
                \begin{scope}[shift={(18,0)}]
                    \node[left, amber, font=\headerfont\footnotesize\bfseries] at (-0.5, 0.75) {VASODILAT.};
                    \draw[thick, amber] (0,0) .. controls (0.2,0) and (0.4,1.5) .. (0.5,1.5) .. controls (0.6,1.5) and (0.7,0.2) .. (0.8,0.2) .. controls (0.9,0.5) and (1.5,0) .. (2,0)
                                         (2,0) .. controls (2.2,0) and (2.4,1.5) .. (2.5,1.5) .. controls (2.6,1.5) and (2.7,0.2) .. (2.8,0.2) .. controls (2.9,0.5) and (3.5,0) .. (4,0);
                \end{scope}
            \end{tikzpicture}
        };
        \fill[cyan, rounded corners=2pt] ([xshift=-0.5pt]box.north west) rectangle ([xshift=3pt]box.south west);
    \end{tikzpicture}\par
\end{minipage}

\vspace{0.4cm}

% --- RODAPÉ ---
\noindent
\begin{tikzpicture}
    \node[rectangle, fill=cardbg, inner sep=10pt, text width=\linewidth-20pt, align=left] {%
        {\headerfont\small\color{cyan}\MakeUppercase{Sumário}} {\small\color{dimtext}$\cdot$ Correlacione a onda com o pulso palpável. Em dúvida, gasometria arterial (SaO\textsubscript{2}) é o padrão-ouro.} \hfill {\headerfont\footnotesize\color{amber}CRITICAL CARE $\cdot$ 2025}
    };
\end{tikzpicture}

\end{minipage}

\end{document}
